\chapter{Data Definitions and Source Provenance}
\label{ch:appendix_a}

This appendix defines the 3DBAG variables, evaluation-set exclusions, and TABULA-NL source decisions needed to reproduce the dataset construction, evaluation, and archetype lookup.
Section~\ref{sec:appendix_a_3dbag} distinguishes the floor count used as a reference attribute from the geometric variables derived for physical calculations and feature analysis.
Section~\ref{sec:appendix_a_evaluation_set} documents the records excluded when the fixed holdout set is reduced to the evaluation set.
Section~\ref{sec:appendix_a_tabula} documents how the operational TABULA-NL lookup was generated and how a difference between the source workbook and the WebTool was resolved.

\section{3DBAG Variables}
\label{sec:appendix_a_3dbag}

3DBAG provides automatically reconstructed geometry rather than surveyed measurements \parencite{threedbagdocs,peters2022,dukai2021}.
Table~\ref{tab:3dbag_fields} separates the reference floor count used in attribute prediction from the geometric variables derived from reconstructed surfaces and heights.
This distinction is necessary because the reference floor count and the floor count used to derive total floor area use different height fields when \texttt{b3\_bouwlagen} is unavailable.

\begin{table}[htbp]
	\centering
	\caption{Definitions and analytical roles of the 3DBAG variables used in this study.}
	\label{tab:3dbag_fields}
	\footnotesize
	\begin{tblr}{
			width = \linewidth,
			colspec = {Q[l,m,3.3cm] X[l] X[l]},
			colsep = 4pt,
			row{1} = {font=\bfseries},
			rowsep = 2pt,
		}
		\toprule
		Variable & Definition & Role in the analysis \\
		\midrule
		Reference floor count & \texttt{b3\_bouwlagen}; if missing, \texttt{round(b3\_h\_50p / 3.0)} & Reference target for floor count prediction \\
		Volume & \texttt{b3\_volume\_lod22} & Reconstructed building volume used in the geometry feature analysis \\
		Envelope area & Sum of exterior wall, flat roof, sloped roof, and ground surface areas; party walls are excluded & Exterior envelope area used to calculate shape factor \\
		Building height & Maximum of zero and \texttt{b3\_h\_max} $-$ \texttt{b3\_h\_maaiveld} & Reconstructed height used when deriving geometric floor count \\
		Geometric floor count & \texttt{b3\_bouwlagen}; if missing, the maximum of one and \texttt{round(building\_height / 3.0)} & Intermediate variable used to derive total floor area \\
		Estimated floor area & \texttt{b3\_opp\_grond} $\times$ \texttt{num\_floors\_estimated} & Floor area used to normalise transmission heat transfer and in the geometry feature analysis \\
		Shape factor & \texttt{envelope\_area / volume} & Exterior surface to volume ratio used in the geometry feature analysis \\
		Shared wall ratio & Party wall area divided by the sum of exterior and party wall areas & Measure of attached wall exposure used in the geometry feature analysis \\
		\bottomrule
	\end{tblr}
\end{table}

The transmission heat transfer calculation uses reconstructed element areas directly and uses \texttt{floor\_area\_estimated} as its denominator.
The geometry feature analysis uses volume, estimated floor area, shape factor, and shared wall ratio as additions to the structured energy class model.
None of these variables is interpreted as a field measurement of the current building envelope.

\section{Evaluation Set Eligibility}
\label{sec:appendix_a_evaluation_set}

The fixed holdout set contains 2,018 buildings, while the performance comparisons use an evaluation set of 2,014 buildings with complete outputs from all three vision configurations and reconstructed areas suitable for the thermal calculation. Reconstructed areas satisfy the validity criteria when estimated floor area is between 20 and 100,000\,m$^2$, exterior wall area is between 1 and 50,000\,m$^2$, and ground area is between 5 and 20,000\,m$^2$. Table~\ref{tab:evaluation_set_exclusions} documents the four excluded buildings. The two exclusion groups do not overlap.

\begin{table}[htbp]
	\centering
	\caption{Buildings excluded from the fixed holdout set when defining the evaluation set.}
	\label{tab:evaluation_set_exclusions}
	\small
	\begin{tblr}{
		width = \linewidth,
		colspec = {Q[l,m,3.8cm] Q[l,m,2.3cm] X[l]},
		colsep = 4pt,
		row{1} = {font=\bfseries},
		rowsep = 3pt,
	}
		\toprule
		\texttt{pand\_id} & Study area & Exclusion reason \\
		\midrule
		\texttt{0599100000240334} & Rotterdam & InternVL3-2B did not return complete predictions for building type, construction year, and floor count \\
		\texttt{0599100000647213} & Rotterdam & InternVL3-2B did not return complete predictions for building type, construction year, and floor count \\
		\texttt{0363100012160228} & Amsterdam & Reconstructed ground and estimated floor areas were below the validity ranges \\
		\texttt{0599100000700053} & Rotterdam & Estimated floor area exceeded the validity range \\
		\bottomrule
	\end{tblr}
\end{table}

\section{TABULA-NL Lookup Provenance}
\label{sec:appendix_a_tabula}

The operational TABULA-NL lookup was generated from the \texttt{Tab.Building.Constr} sheet of the TABULA values workbook and checked against the Dutch exemplary building records in the TABULA WebTool \parencite{loga2016,tabulawebtool}.
The generator selects existing state construction assemblies for each residential size class and construction period.
Terraced houses use the solid wall variants, while single family houses, multifamily houses, and apartment blocks use the cavity wall variants represented by the corresponding exemplary buildings.
When the workbook contains several subperiod assemblies within one TABULA-NL period, the implementation selects the assembly with the highest U-value. Window U-values are assigned directly by construction period.

This value is used consistently in Table~\ref{tab:tabula_uvalues} and in all archetype and transmission heat transfer calculations.
The generated lookup and its element source mapping are stored in \texttt{data/processed/tabula\_nl.csv} and \texttt{data/processed/tabula\_nl\_provenance.csv}.
